% ============================================================
% 纯答案册·课时13-2.6.1双曲线的标准方程 —— 工具/答案抽册器.py 抽册生成件（勿手改；第三档＝纯答案单独成册）
% 源件（只读）：C:/提示词/工作区/M3-第2章量产0913/成卷/导学件/课时13-2.6.1双曲线的标准方程/main.tex（md5 5b9bec74a29bf71d51140bd1254c68cb）
%% 口径：题面不重印；逐题＝题号(印面号同正册)＋[答案]＋[详解]，值/详解照源件逐字
%       （钉值门硬断言：册值≡正册件内值≡值台账值tex，逐字全等）。册序＝正册装配序＝印面号 1..21 升序（同序同号）；键序断言基准＝题面库冻结 manifest canonical 键序。
%% 三档导出：整体 main-true／纯题 main-false（在产）｜本册＝纯答案册（本件）
%% 双档：ansbook-true.tex＝详解印本；ansbook-false.tex＝纯值速查本（详解不印）
% 编译：xelatex <壳> ×2，cwd＝本目录；sty＝qp-m3.sty 本地挂载（md5 7c3930362be8a0a2bdf21bbf8ac16573，锁校验）
% ============================================================
\documentclass[fontset=none]{ctexart}
\usepackage{qp-m3}
% —— 印面逐字符号路由补钉（承源件；− 走 TNR、▱ NSC 子块；禁改 sty 保 md5） ——
\xeCJKDeclareCharClass{Default}{"2212}%
\setCJKfamilyfont{nscblk}[Path=C:/提示词/工作区/字替对照-0909/variantF/fonts/,BoldFont=NSC-w600.ttf]{NSC-w500.ttf}%
\xeCJKDeclareSubCJKBlock{nscblk}{"25B1}%
\newcommand{\pxparallelogram}{{\CJKfamily{nscblk}▱}}%
% —— 断行微调（承母版实测档，三零门承重；\emergencystretch 不另设） ——
\xeCJKsetup{CJKglue={\hskip 0.10em plus 0.08em minus 0.05em}}%
% —— 纯答案册全线括线（长详解可跨栏断，承练习件全线括线口径；灰底盒不可跨栏断） ——
\ansblockgrayfalse
% —— 双档开关（false 档＝纯值速查：答案印、详解不印；件面开关，不动 sty 本体） ——
\newif\ifansbookdetail
\ifdefined\ansbookdetail
  \ifnum\ansbookdetail=0 \ansbookdetailfalse\else\ansbookdetailtrue\fi
\else
  \ansbookdetailtrue
\fi
% —— 页脚件名（承源件章名；件名＝<件型>答案册） ——
\renewcommand{\qpzhangming}{第二章\quad 平面解析几何}%
\renewcommand{\qpjianming}{导学件答案册}%

\begin{document}
\zhangtitle{第二章\quad 平面解析几何}
\jietitle{2.6.1 双曲线的标准方程}
\par\glueguard{1}\addvspace{2pt}\noindent\biaoqian{【纯答案册】}{\fangsong 题面不重印\quad 印面号与正册同号同序}\par\addvspace{2pt}

\begin{multicols}{2}
\emergencystretch=1em

\begin{ansblock}[2章-练-课时13-简1]
% ans:2章-练-课时13-简1
\ansitem{1}{(1)存在，轨迹是线段 AB；(2)存在，轨迹是以 B 为端点、指向 A 一侧反向（即线段 AB 的延长线方向）的射线；(3)存在，轨迹是以 A 为端点、沿线段 BA 延长线的射线；(4)不存在；(5)不存在}
\ifansbookdetail
\ansnote{详解}{记\(|AB|=6\)，对动点\(P\)总有\(|PA|+|PB|\geqslant|AB|\)（三角不等式）与\(\big||PA|-|PB|\big|\leqslant|AB|\)（两边之差小于等于第三边，等号当\(P\)、\(A\)、\(B\)共线）．(1)\(|PA|+|PB|=6=|AB|\)：等号成立当且仅当\(P\)在线段\(AB\)上，故轨迹为线段\(AB\)（椭圆定义中\(2a=2c\)的退化位）．(2)\(|PA|-|PB|=6=|AB|\)：须\(\big||PA|-|PB|\big|=|AB|\)，即\(P\)、\(A\)、\(B\)共线且\(P\)不在\(A\)、\(B\)之间；再由\(|PA|>|PB|\)知\(P\)靠\(B\)一侧，即\(P\)在射线（以\(B\)为端点、方向背离\(A\)）上，此即“线段\(AB\)的延长线”．(3)\(|PB|-|PA|=6\)：同(2)对称，\(P\)在以\(A\)为端点、背离\(B\)的射线上（“线段\(BA\)的延长线”）．(4)\(|PA|+|PB|=3<6=|AB|\)：与三角不等式矛盾，满足条件的\(P\)不存在（轨迹不存在）．(5)\(|PB|-|PA|=7>6=|AB|\)：与\(\big||PB|-|PA|\big|\leqslant|AB|\)矛盾，\(P\)不存在（轨迹不存在）．}
\fi
\end{ansblock}

\begin{ansblock}[2章-练-课时13-简9]
% ans:2章-练-课时13-简9
\ansitem{2}{x²/4−y²/5＝1}
\ifansbookdetail
\ansnote{详解}{\(|F_1F_2|=6>4>0\)，由双曲线定义：\(2a=4\)，\(a=2\)；\(c=3\)；\(b^{2}=c^{2}-a^{2}=9-4=5\)．焦点在\(x\)轴，方程\(x^{2}/4-y^{2}/5=1\)．辨析：若差为\(6=|F_1F_2|\)，轨迹为两条射线\(y=0\)（\(x\leqslant -3\)或\(x\geqslant 3\)）；若差大于\(6\)，则无轨迹．}
\fi
\end{ansblock}

\begin{ansblock}[2章-练-课时13-简3]
% ans:2章-练-课时13-简3
\ansitem{3}{(1)k＜−3 或 1＜k＜3；(2)1＜k＜3；(3)k＜−3}
\ifansbookdetail
\ansnote{详解}{两边乘\(-1\)化为标准形：\(x^{2}/(k-1)+y^{2}/(|k|-3)=1\)（须\(k-1\neq 0\)且\(|k|-3\neq 0\)）．(1)表示双曲线当且仅当两项系数异号，即\((k-1)(|k|-3)<0\)：\(k-1>0\)且\(|k|-3<0\)，即\(k>1\)且\(-3<k<3\)，得\(1<k<3\)；\(k-1<0\)且\(|k|-3>0\)，即\(k<1\)且（\(k>3\)或\(k<-3\)），得\(k<-3\)；故\(k<-3\)或\(1<k<3\)．(2)焦点在\(x\)轴上的双曲线当且仅当\(x^{2}\)项系数为正、\(y^{2}\)项为负：\(k-1>0\)且\(|k|-3<0\)，得\(1<k<3\)．(3)焦点在\(y\)轴上的双曲线当且仅当\(y^{2}\)项系数为正、\(x^{2}\)项为负：\(|k|-3>0\)且\(k-1<0\)，得\(k<-3\)（\(k>3\)与\(k<1\)无交）．}
\fi
\end{ansblock}

\begin{ansblock}[2章-练-课时13-简4]
% ans:2章-练-课时13-简4
\ansitem{4}{(1)x²/6 − y²/8 ＝1；(2)y²/9 − x²/16 ＝1；(3)x²/8 − y²/8 ＝1；(4)y²/16 − x²/9 ＝1}
\ifansbookdetail
\ansnote{详解}{(1)点\((\sqrt{6},0)\)在\(x\)轴上且在曲线上，故焦点在\(x\)轴、该点即右顶点，\(a^{2}=6\)；设\(x^{2}/6-y^{2}/b^{2}=1\)，代\((3,2)\)：\(9/6-4/b^{2}=1\)，即\(4/b^{2}=1/2\)，\(b^{2}=8\)，得\(x^{2}/6-y^{2}/8=1\)（核验\(9/6-4/8=3/2-1/2=1\)成立）．(2)焦点在\(y\)轴、\(c=5\)．设\(P(4\sqrt{3}/3,2\sqrt{3})\)，到两焦点\((0,\pm 5)\)的距离平方：\(|PF_1|^{2}=(4\sqrt{3}/3)^{2}+(2\sqrt{3}-5)^{2}=16/3+37-20\sqrt{3}=(10-3\sqrt{3})^{2}/3\)，故\(|PF_1|=10\sqrt{3}/3-3\)；\(|PF_2|^{2}=(10+3\sqrt{3})^{2}/3\)，故\(|PF_2|=10\sqrt{3}/3+3\)；于是\(2a=\big||PF_2|-|PF_1|\big|=6\)，即\(a=3\)、\(b^{2}=c^{2}-a^{2}=25-9=16\)，得\(y^{2}/9-x^{2}/16=1\)；回代\(P\)：\(12/9-(16/3)/16=4/3-1/3=1\)成立（焦点在\(y\)轴成立）．(3)\(a=b\)为等轴双曲线，两轴取向皆可能，设\(x^{2}-y^{2}=t\)（\(t\neq 0\)），代\((3,-1)\)：\(t=9-1=8>0\)，得\(x^{2}-y^{2}=8\)，即\(x^{2}/8-y^{2}/8=1\)．(4)焦点位置未定，设\(mx^{2}-ny^{2}=1\)（\(mn>0\)），代入两点：\(9m-32n=1\)、\((81/16)m-25n=1\)，由第一式\(m=(1+32n)/9\)，代第二式\((9/16)(1+32n)-25n=1\)，即\(9/16+18n-25n=1\)，\(-7n=7/16\)，\(n=-1/16\)、\(m=-1/9\)，方程为\(-x^{2}/9+y^{2}/16=1\)，即\(y^{2}/16-x^{2}/9=1\)（回代\((9/4,5)\)：\(25/16-(81/16)/9=25/16-9/16=1\)成立）．}
\fi
\end{ansblock}

\begin{ansblock}[2章-练-课时13-简10]
% ans:2章-练-课时13-简10
\ansitem{5}{x²/9−y²/16＝1；焦点 (±5,0)；e=5/3}
\ifansbookdetail
\ansnote{详解}{\(2a=6\)即\(a=3\)；\(2b=8\)即\(b=4\)；\(c^{2}=a^{2}+b^{2}=9+16=25\)，\(c=5\)．故标准方程\(x^{2}/9-y^{2}/16=1\)，焦点\((\pm 5,0)\)，\(e=c/a=5/3\)．}
\fi
\end{ansblock}

\begin{ansblock}[2章-练-课时13-简2]
% ans:2章-练-课时13-简2
\ansitem{6}{C}
\ifansbookdetail
\ansnote{详解}{\(a^{2}=25\)、\(b^{2}=23\)，即\(a=5\)、\(2a=10\)，\(c^{2}=48\)、\(c=4\sqrt{3}\)．由定义\(\big||PF_2|-|PF_1|\big|=2a=10\)，即\(\big||PF_2|-8\big|=10\)，得\(|PF_2|=18\)或\(|PF_2|=-2\)（距离非负，舍），故\(|PF_2|=18\)．（存在性核验：\(c-a=4\sqrt{3}-5\approx 1.93\)、\(c+a\approx 11.93\)；\(|PF_1|=8\in[c-a,c+a)\)只能对应\(P\)在左支，此时\(|PF_2|=|PF_1|+2a=18\geqslant c+a\)成立，与选项A、B中“2”所设的右支情形（须\(|PF_2|\geqslant c-a\)且\(|PF_1|\geqslant c+a\)）不合，故不取双解．）故选：C．}
\fi
\end{ansblock}

\begin{ansblock}[2章-练-课时13-简5]
% ans:2章-练-课时13-简5
\ansitem{7}{C}
\ifansbookdetail
\ansnote{详解}{\(A\)、\(B\)在同一支上，由定义\(|AF_2|-|AF_1|=2a\)、\(|BF_2|-|BF_1|=2a\)（该支上的点到另一焦点较远，故取此号），两式相加：\((|AF_2|+|BF_2|)-(|AF_1|+|BF_1|)=4a\)．又弦\(AB\)过焦点\(F_1\)，\(F_1\)在\(A\)、\(B\)之间（焦点在该支围成的开口内），故\(|AF_1|+|BF_1|=|AB|=m\)，于是\(|AF_2|+|BF_2|=4a+m\)，\(\triangle ABF_2\)的周长\(=|AF_2|+|BF_2|+|AB|=4a+m+m=4a+2m\)．故选：C．}
\fi
\end{ansblock}

\begin{ansblock}[2章-练-课时13-简6]
% ans:2章-练-课时13-简6
\ansitem{8}{C}
\ifansbookdetail
\ansnote{详解}{\(a^{2}=9\)、\(b^{2}=16\)，即\(a=3\)、\(b=4\)，\(c^{2}=a^{2}+b^{2}=25\)、\(c=5\)，故左焦点\(F(-5,0)\)，虚轴长\(2b=8\)，\(|PQ|=2\times 8=16\)．点\(A(5,0)\)正是右焦点\(F'\)，且\(A\)在线段\(PQ\)上，故\(|PF|-|PA|=|PF|-|PF'|=2a=6\)（\(P\)在右支），同理\(|QF|-|QA|=6\)，两式相加：\(|PF|+|QF|=12+(|PA|+|QA|)=12+|PQ|=12+16=28\)，所以\(\triangle PFQ\)的周长\(=|PF|+|QF|+|PQ|=28+16=44\)．故选：C．}
\fi
\end{ansblock}

\begin{ansblock}[2章-练-课时13-简7]
% ans:2章-练-课时13-简7
\ansitem{9}{B}
\ifansbookdetail
\ansnote{详解}{\(a^{2}=3-m\)、\(b^{2}=m\)（\(0<m<3\)保证两项均正、焦点在\(x\)轴），\(c^{2}=a^{2}+b^{2}=3\)，即\(c=\sqrt{3}\)、\(|F_1F_2|=2\sqrt{3}\)．\(|OP|=\frac{1}{2}|F_1F_2|=\sqrt{3}=|OF_1|=|OF_2|\)，即\(P\)到线段\(F_1F_2\)中点\(O\)的距离等于该线段的一半，故\(P\)在以\(F_1F_2\)为直径的圆上，\(\angle F_1PF_2=90^{\circ}\)（直径所对圆周角为直角）．在直角\(\triangle F_1PF_2\)中，\(\angle PF_1F_2=30^{\circ}\)，故\(|PF_2|=|F_1F_2|\sin 30^{\circ}=2\sqrt{3}\times\frac{1}{2}=\sqrt{3}\)，\(|PF_1|=|F_1F_2|\cos 30^{\circ}=2\sqrt{3}\times\frac{\sqrt{3}}{2}=3\)．\(P\)在右支上，由定义\(|PF_1|-|PF_2|=2a\)：\(3-\sqrt{3}=2\sqrt{3-m}\)，即\(\sqrt{3-m}=(3-\sqrt{3})/2\)，\(3-m=(9-6\sqrt{3}+3)/4=3-\frac{3\sqrt{3}}{2}\)，得\(m=\frac{3\sqrt{3}}{2}\)．（核验\(0<m<3\)：\(3\sqrt{3}/2\approx 2.598<3\)成立；此时\(a^{2}=3-m\approx 0.402>0\)成立，\(2a=3-\sqrt{3}=|PF_1|-|PF_2|\)成立．）故选：B．}
\fi
\end{ansblock}

\begin{ansblock}[2章-练-课时13-简8]
% ans:2章-练-课时13-简8
\ansitem{10}{A}
\ifansbookdetail
\ansnote{详解}{\(a^{2}=4\)、\(b^{2}=12\)，即\(a=2\)、\(c^{2}=16\)、\(c=4\)，故左焦点\(F(-4,0)\)、右焦点\(F'(4,0)\)．\(P\)在右支上，由定义\(|PF|-|PF'|=2a=4\)，即\(|PF|=|PF'|+4\)，于是\(|PF|+|PA|=|PF'|+|PA|+4\geqslant|AF'|+4\)（三角不等式），\(|AF'|=\sqrt{(1-4)^{2}+4^{2}}=\sqrt{25}=5\)，故\(|PF|+|PA|\geqslant 5+4=9\)．取等核验：须\(P\)在线段\(AF'\)上，联立直线\(AF'\)（过\(A(1,4)\)、\(F'(4,0)\)，方程\(y=(16-4x)/3\)）与右支\(x^{2}/4-y^{2}/12=1\)（\(x\geqslant 2\)）得\(11x^{2}+128x-364=0\)，解得\(x=26/11\)（另一根\(-14\)舍），\(26/11\approx 2.36\in[2,4]\)，对应\(y=24/11\)，交点\((26/11,24/11)\)确实落在线段\(AF'\)上且在右支上，故等号可取，最小值为9．故选：A．}
\fi
\end{ansblock}

\begin{ansblock}[2章-练-课时13-中1]
% ans:2章-练-课时13-中1
\ansitem{11}{4√3}
\ifansbookdetail
\ansnote{详解}{焦点\((\pm 4,0)\)即\(c=4\)、\(c^{2}=16\)；又\(c^{2}=m+4\)，即\(m=12\)，故\(a^{2}=12\)、\(a=2\sqrt{3}\)、\(2a=4\sqrt{3}\)．设\(|MF_1|=t_1\)、\(|MF_2|=t_2\)，由定义\(|t_1-t_2|=2a=4\sqrt{3}\)，平方得\(t_1^{2}+t_2^{2}-2t_1t_2=48\)……①在\(\triangle F_1MF_2\)中用余弦定理：\(t_1^{2}+t_2^{2}-2t_1t_2\cos 60^{\circ}=|F_1F_2|^{2}=64\)，即\(t_1^{2}+t_2^{2}-t_1t_2=64\)……②②\(-\)①得\(t_1t_2=64-48=16\)，所以\(S=\frac{1}{2}t_1t_2\sin 60^{\circ}=\frac{1}{2}\times 16\times\frac{\sqrt{3}}{2}=4\sqrt{3}\)．}
\fi
\end{ansblock}

\begin{ansblock}[2章-练-课时13-中2]
% ans:2章-练-课时13-中2
\ansitem{12}{16/5}
\ifansbookdetail
\ansnote{详解}{\(a^{2}=9\)、\(b^{2}=16\)，即\(c^{2}=25\)、\(c=5\)、\(2a=6\)、\(|F_1F_2|=10\)．设\(|PF_1|=m\)、\(|PF_2|=n\)，由定义\(\big||PF_1|-|PF_2|\big|=2a=6\)，勾股（\(PF_1\perp PF_2\)）\(m^{2}+n^{2}=|F_1F_2|^{2}=100\)，两式并立：\((m-n)^{2}=m^{2}+n^{2}-2mn\)，即\(36=100-2mn\)，得\(mn=32\)．\(\triangle PF_1F_2\)的面积既等于\(\frac{1}{2}mn\)（直角边互乘），又等于\(\frac{1}{2}|F_1F_2|\cdot h\)（以焦轴为底、\(h\)为\(P\)到\(x\)轴的距离）：\(\frac{1}{2}\times 32=\frac{1}{2}\times 10\times h\)，即\(16=5h\)，得\(h=16/5\)．（存在性核验：由\(m-n=\pm 6\)、\(mn=32\)得\((m+n)^{2}=100+64=164\)，\(m+n=2\sqrt{41}\)，故\(\{m,n\}=\{\sqrt{41}+3,\sqrt{41}-3\}\approx\{9.4,3.4\}\)，两值均正且平方和为\(2(41+9)=100\)成立；与双曲线上焦半径的取值范围（近焦点侧\(\geqslant c-a=2\)、远焦点侧\(\geqslant c+a=8\)）比对：\(9.4\geqslant 8\)配远侧、\(3.4\geqslant 2\)配近侧成立，故满足垂直条件的点\(P\)确实存在．）所以点\(P\)到\(x\)轴的距离为\(16/5\)．}
\fi
\end{ansblock}

\begin{ansblock}[2章-练-课时13-中3]
% ans:2章-练-课时13-中3
\ansitem{13}{D}
\ifansbookdetail
\ansnote{详解}{\(a^{2}=16\)、\(b^{2}=9\)，即\(c^{2}=25\)、\(c=5\)，故\(A(-5,0)\)正是左焦点，记右焦点\(F(5,0)\)．\(P\)在右支上，由定义\(|PA|-|PF|=2a=8\)，即\(|PA|=|PF|+8\)，于是\(|PA|-|PB|=8+(|PF|-|PB|)\leqslant 8+|BF|\)（三角形两边之差小于第三边），\(|BF|=\sqrt{(8-5)^{2}+(4-0)^{2}}=\sqrt{25}=5\)，故\(|PA|-|PB|\leqslant 13\)．取等核验：须\(B\)在线段\(PF\)上，即\(P\)在射线\(F\to B\)的延长方向上；该射线参数式\((5+3s,4s)\)（\(s>0\)时经\(B\)，对应\(s=1\)），代入\(9x^{2}-16y^{2}=144\)得\(175s^{2}-270s-81=0\)，正根\(s=9/5>1\)，对应\(P(52/5,36/5)\)（核验\(9\times(52/5)^{2}-16\times(36/5)^{2}=(24336-20736)/25=144\)成立，且\(x=52/5\geqslant 4\)在右支上成立），故等号可取，最大值为13．故选：D．}
\fi
\end{ansblock}

\begin{ansblock}[2章-练-课时13-中4]
% ans:2章-练-课时13-中4
\ansitem{14}{x²/4 − y²/5 ＝1}
\ifansbookdetail
\ansnote{详解}{记\(F_1(-3,0)\)、\(F_2(3,0)\)，两定圆半径均为2，\(|F_1F_2|=6\)．设动圆\(M\)半径为\(r\)．情形一（与圆\(F_1\)内切、与圆\(F_2\)外切）：\(|MF_1|=r-2\)或\(2-r\)；\(|MF_2|=r+2\)．若\(|MF_1|=2-r\)，则\(|MF_1|+|MF_2|=4<|F_1F_2|=6\)，不可能（三角不等式），故\(|MF_1|=r-2\)（须\(r\geqslant 2\)），此时\(|MF_2|-|MF_1|=4\)；情形二（与圆\(F_1\)外切、与圆\(F_2\)内切）：同理\(|MF_1|=r+2\)、\(|MF_2|=r-2\)，\(|MF_1|-|MF_2|=4\)；（动圆含入定圆内的情形如上所验均导致\(|MF_1|+|MF_2|=4<6\)，故皆不出现）合写为\(\big||MF_1|-|MF_2|\big|=4<|F_1F_2|=6\)，由双曲线定义，\(M\)的轨迹是以\(F_1\)、\(F_2\)为焦点的双曲线（情形一给左支、情形二给右支，两支并存），\(2a=4\)、\(2c=6\)，即\(a=2\)、\(c=3\)、\(b^{2}=c^{2}-a^{2}=9-4=5\)，轨迹方程为\(x^{2}/4-y^{2}/5=1\)．}
\fi
\end{ansblock}

\begin{ansblock}[2章-练-课时13-难1]
% ans:2章-练-课时13-难1
\ansitem{15}{9/2}
\ifansbookdetail
\ansnote{详解}{设公共半焦距为\(c\)，椭圆长半轴\(a_1\)、双曲线实半轴\(a_2\)，记\(m=|PF_1|\)、\(n=|PF_2|\)．由椭圆定义\(m+n=2a_1\)、双曲线定义\(|m-n|=2a_2\)，又\(\angle F_1PF_2=90^{\circ}\)得\(m^{2}+n^{2}=4c^{2}\)；将前两式分别平方：\((m+n)^{2}+(m-n)^{2}=2(m^{2}+n^{2})\)，即\(4a_1^{2}+4a_2^{2}=2\times 4c^{2}\)，故\(a_1^{2}+a_2^{2}=2c^{2}\)，两边除\(c^{2}\)：\(1/e_1^{2}+1/e_2^{2}=2\)……①于是\(4e_1^{2}+e_2^{2}=\frac{1}{2}(1/e_1^{2}+1/e_2^{2})(4e_1^{2}+e_2^{2})=\frac{1}{2}\bigl[4+e_2^{2}/e_1^{2}+4e_1^{2}/e_2^{2}+1\bigr]=\frac{1}{2}\bigl[5+e_2^{2}/e_1^{2}+4e_1^{2}/e_2^{2}\bigr]\geqslant\frac{1}{2}\bigl[5+2\sqrt{(e_2^{2}/e_1^{2})\cdot(4e_1^{2}/e_2^{2})}\bigr]=\frac{1}{2}(5+4)=9/2\)（基本不等式）．取等条件：\(e_2^{2}/e_1^{2}=4e_1^{2}/e_2^{2}\)，即\(e_2^{2}=2e_1^{2}\)，代①：\(1/e_1^{2}+1/(2e_1^{2})=2\)，即\(3/(2e_1^{2})=2\)，得\(e_1^{2}=3/4\)、\(e_2^{2}=3/2\)，即\(e_1=\sqrt{3}/2\in(0,1)\)成立（椭圆）、\(e_2=\sqrt{6}/2\in(1,\sqrt{2})\)成立（双曲线），且\(a_1=c/e_1=2c/\sqrt{3}>c\)、\(a_2=c/e_2=\sqrt{2/3}\,c<c\)均合要求；故\(4e_1^{2}+e_2^{2}\)的最小值为\(9/2\)．}
\fi
\end{ansblock}

\begin{ansblock}[2章-练-课时13-难2]
% ans:2章-练-课时13-难2
\ansitem{16}{(1)(4,0)；(2)没有"被抓"的风险}
\ifansbookdetail
\ansnote{详解}{(1)设机器鼠位置为\(P\)，声音传播时间差为距离差除以\(v_0\)：\(|PA|/v_0-|PB|/v_0=8/v_0\)，即\(|PA|-|PB|=8\)（米），而\(|AB|=10\)，\(8<10\)，由双曲线定义\(P\)的轨迹是以\(A\)、\(B\)为焦点的双曲线中靠\(B\)的一支（右支）：\(2a=8\)、\(2c=10\)，即\(a=4\)、\(c=5\)、\(b^{2}=c^{2}-a^{2}=9\)，轨迹方程\(x^{2}/16-y^{2}/9=1\)（\(x\geqslant 4\)）．时刻\(t_0\)有\(|OP|=4\)：对右支上的点\(x\geqslant 4\)，故\(|OP|=\sqrt{x^{2}+y^{2}}\geqslant x\geqslant 4\)，等号仅当\(x=4\)且\(y=0\)，即顶点\((4,0)\)，所以\(t_0\)时机器鼠所在位置坐标为\((4,0)\)．(2)直线\(l\)：\(y=x\)，设与其平行的直线\(l_1\)：\(y=x+m\)与右支相切．代入\(9x^{2}-16y^{2}=144\)：\(9x^{2}-16(x+m)^{2}=144\)，即\(7x^{2}+32mx+16m^{2}+144=0\)，\(\Delta=(32m)^{2}-4\times 7\times(16m^{2}+144)=1024m^{2}-448m^{2}-4032=576m^{2}-4032=0\)，得\(m^{2}=7\)，切点横坐标\(x=-32m/(2\times 7)=-16m/7\)，要在右支须\(x\geqslant 4\)，故取\(m=-\sqrt{7}\)（此时\(x=16\sqrt{7}/7\approx 6.05\geqslant 4\)成立；\(m=+\sqrt{7}\)的切点在左支，舍），即与右支相切且最“靠向”\(l\)的平行线为\(l_1\)：\(y=x-\sqrt{7}\)，切点\(T(16\sqrt{7}/7,9\sqrt{7}/7)\)即右支上到\(l\)最近的点．两平行线\(y=x\)与\(y=x-\sqrt{7}\)的距离\(d=|\sqrt{7}|/\sqrt{1+1}=\sqrt{14}/2\approx 1.87\)（米）\(>1.5\)，（且当\(|m|<\sqrt{7}\)时\(\Delta<0\)，即形如\(y=x+m\)的直线与双曲线无公共点，右支整体落在\(x-y\geqslant\sqrt{7}\)的一侧，故此距离确为最小值而非极值假象）所以机器鼠到直线\(l\)的最小距离为\(\sqrt{14}/2\approx 1.87\)米，大于1.5米，没有“被抓”的风险．}
\fi
\end{ansblock}

\begin{ansblock}[2章-导-课时13-G1]
% ans:2章-导-课时13-G1
\ansitem{17}{C}
\ifansbookdetail
\ansnote{详解}{\(a^{2}=16\)、\(b^{2}=48\)，即\(a=4\)，\(c^{2}=a^{2}+b^{2}=64\)、\(c=8\)，故\(a+c=12\)、\(c-a=4\)．先定支：若\(P\)在右支上，则\(|PF_1|\geqslant a+c=12\)（右支上的点到左焦点的距离以右顶点处最小，为\(a+c\)），与\(|PF_1|=10\)矛盾；所以\(P\)在左支上，左支上\(|PF_2|>|PF_1|\)，由定义\(|PF_2|-|PF_1|=2a=8\)，得\(|PF_2|=10+8=18\)．（另一解2的来源是误按“\(|10-|PF_2||=8\)”取\(|PF_2|=2\)，但此时\(P\)须在右支且\(|PF_2|\geqslant c-a=4\)，不合，舍．）故选：C．}
\fi
\end{ansblock}

\begin{ansblock}[2章-导-课时13-G2]
% ans:2章-导-课时13-G2
\ansitem{18}{D}
\ifansbookdetail
\ansnote{详解}{设动圆\(P\)的半径为\(r\)，则\(|PM|=r\)；圆\(N\)的圆心\(N(4,0)\)、半径4，且\(|MN|=8\)．两圆相切分三种位置：①外切\(|PN|=r+4\)；②圆\(N\)在动圆内（内切）\(|PN|=r-4\)（须\(r\geqslant 4\)）；③动圆在圆\(N\)内（内切）\(|PN|=4-r\)．情形③不可能：此时\(|PM|+|PN|=r+4-r=4<|MN|=8\)，与三角形两边之和不小于第三边矛盾（也即\(M\)在圆\(N\)外，动圆既然过\(M\)就不能整个含在圆\(N\)内）．情形①：\(|PN|-|PM|=4\)；情形②：\(|PM|-|PN|=4\)．合写即\(\big||PN|-|PM|\big|=4<|MN|=8\)，由双曲线定义，\(P\)的轨迹是以\(M\)、\(N\)为焦点的双曲线，\(2a=4\)、\(2c=8\)，即\(a=2\)、\(c=4\)、\(b^{2}=c^{2}-a^{2}=12\)，且两种切法都能实现（在双曲线上任取一点\(P\)，取\(r=|PM|\)即可作出满足条件的动圆），故两支并存、不加\(x\)的限制，轨迹方程为\(x^{2}/4-y^{2}/12=1\)．故选：D．}
\fi
\end{ansblock}

\begin{ansblock}[2章-导-课时13-G3]
% ans:2章-导-课时13-G3
\ansitem{19}{B}
\ifansbookdetail
\ansnote{详解}{\(N\)是\(F_1M\)的中点（对称关系）、\(O\)是\(F_1F_2\)的中点，故\(ON\)是\(\triangle MF_1F_2\)的中位线，即\(|MF_2|=2|ON|=2\)（\(N\)在单位圆上）．\(P\)在\(F_1M\)的垂直平分线上，即\(|PF_1|=|PM|\)；又\(P\)在直线\(F_2M\)上，故\(P\)、\(M\)、\(F_2\)三点共线，于是\(\big||PF_2|-|PF_1|\big|=\big||PF_2|-|PM|\big|=|MF_2|=2<|F_1F_2|=4\)，由双曲线定义，\(P\)的轨迹是以\(F_1\)、\(F_2\)为焦点的双曲线（\(2a=2\)、\(2c=4\)，即\(x^{2}-y^{2}/3=1\)）．故选：B．}
\fi
\end{ansblock}

\begin{ansblock}[2章-导-课时13-G4]
% ans:2章-导-课时13-G4
\ansitem{20}{√2}
\ifansbookdetail
\ansnote{详解}{把三点代入\(x^{2}/a^{2}-y^{2}\)的值：\((-2,1)\)与\((2,-1)\)都得\(4/a^{2}-1\)，\((-2,3)\)得\(4/a^{2}-9\)．先试\((-2,3)\)在\(C\)上：\(4/a^{2}-9=1\)，即\(a^{2}=2/5\)，此时另两点处的值\(4/a^{2}-1=10-1=9\neq 1\)，即只有1个点在\(C\)上，与“恰有两个”矛盾，故\((-2,3)\)不在\(C\)上．于是恰有的两个点只能是\((-2,1)\)与\((2,-1)\)（两点关于原点对称，双曲线关于原点中心对称，故同属\(C\)成立），由\(4/a^{2}-1=1\)，即\(a^{2}=2\)，得\(a=\sqrt{2}\)（\(a>0\)）．（回代核验：\(a^{2}=2\)时\(4/2-1=1\)成立，两点在\(C\)上；\(4/2-9=-7\neq 1\)成立，第三点不在，恰两个成立．）}
\fi
\end{ansblock}

\begin{ansblock}[2章-导-课时13-G5]
% ans:2章-导-课时13-G5
\ansitem{21}{x²/9 − y²/16 ＝1 (x≠±3)}
\ifansbookdetail
\ansnote{详解}{设\(A(x,y)\)，则两斜率存在须\(x\neq\pm 3\)（\(x=\pm 3\)时\(A\)与\(B\)或\(C\)的横坐标相同，直线\(AB\)或\(AC\)垂直于\(x\)轴，斜率不存在，且\(A\)与\(B\)或\(C\)重合不成三角形），\(k_{AB}=y/(x+3)\)、\(k_{AC}=y/(x-3)\)，由题意\(y^{2}/[(x+3)(x-3)]=16/9\)，即\(y^{2}/(x^{2}-9)=16/9\)，去分母：\(9y^{2}=16(x^{2}-9)\)，即\(16x^{2}-9y^{2}=144\)，得\(x^{2}/9-y^{2}/16=1\)．该双曲线上\(x^{2}\geqslant 9\)，等号只在顶点\((\pm 3,0)\)，而这两点使\(A\)与\(B\)、\(C\)重合（三角形退化），故除去\(x=\pm 3\)两点，即\(x\neq\pm 3\)（此时自动满足\(|x|>3\)，两斜率均存在成立）．所以顶点\(A\)的轨迹方程为\(x^{2}/9-y^{2}/16=1\)（\(x\neq\pm 3\)）．}
\fi
\end{ansblock}

% ---- 尾块（\tailfill[30mm] 定高参——钉档 §三.3：无参在平衡栏呈「内容尾随框」，贴栏底须定高参；实测无参致末页平衡 Overfull\vbox 12.99pt，定高 30mm 三零） ----
\tailfill[30mm]

\end{multicols}
\end{document}
